\section*{Capítulo \ref{ch:b1:erosion-transport} --- Erosão e transporte}

\begin{solution}{exo:b1:erosion-transport:1}
\qty{0.4}{g.L^{-1}} é \qty{0.4}{kg.m^{-3}}, logo o fluxo é $500 \times 0.4 =
\qty{200}{kg.s^{-1}}$. Um ano de \num{365.25}~dias dura \qty{3.156e7}{s}, logo o fluxo anual é \qty{6.31e9}{kg}, ou \qty{6.31}{Mt}.
\end{solution}

\begin{solution}{exo:b1:erosion-transport:2}
$w_s = R g d^2/(18\nu) = 1.6548 \times 9.81 \times (\num{30e-6})^2/(18 \times
\num{1.0034e-6}) = \qty{8.09e-4}{m.s^{-1}}$, cerca de \qty{0.8}{mm} por segundo. Seu número de Reynolds do grão é $w_s d/\nu = \num{8.09e-4} \times
\num{30e-6}/\num{1.0034e-6} = 0.024$, muito abaixo de 1: vale a lei de Stokes. Ele afunda através de \qty{2}{m} em $2/\num{8.09e-4} = \qty{2472}{s}$, cerca de
\qty{41}{min}.
\end{solution}

\begin{solution}{exo:b1:erosion-transport:3}
$\tau_b = \rho g h S = 998.2 \times 9.81 \times 0.5 \times \num{4e-4} =
\qty{1.96}{Pa}$. Com $(\rho_s - \rho) g = 1651.8 \times 9.81 =
\qty{16204}{N.m^{-3}}$: para \qty{1}{mm}, $\theta = 1.96/(16204 \times 0.001) =
0.121$, quase quatro vezes o limiar de \num{0.0314}, logo a areia se move; para
\qty{10}{mm}, $\theta = 0.0121$, um quinto do limiar de \num{0.0556}, logo o cascalho permanece. O riacho carrega sua areia sobre um leito de cascalho que não se move.
\end{solution}

\begin{solution}{exo:b1:erosion-transport:4}
(a)~$12.6/3.9 = 3.2$: os rios carregam cerca de três vezes mais em suspensão que em solução. (b)~Grandes reservatórios: $0.20 \times 14 = \qty{2.8}{Gt}$ por ano; todos eles: $0.26 \times 14 = \qty{3.64}{Gt}$ por ano. (c)~O fluxo dissolvido é de \qty{3.9e15}{g} por ano a \qty{120}{g.m^{-3}}, logo a água que o carrega é $\num{3.9e15}/120 = \qty{3.25e13}{m^3}$, cerca de \qty{32500}{km^3} por ano. A \omterm{def:b1:erosion-transport:load}{carga em suspensão} nessa água é $\qty{1.26e16}{g}/\qty{3.25e16}{L} =
\qty{0.39}{g.L^{-1}}$.
\end{solution}

\begin{solution}{exo:b1:erosion-transport:5}
(a)~A seção tem a área $10 \times 2 = \qty{20}{m^2}$ e o perímetro molhado $10 + 2 \times 2 = \qty{14}{m}$, logo $R_h = 20/14 = \qty{1.43}{m}$. (b)~$\tau_b = 998.2 \times 9.81 \times 1.43 \times 0.001 = \qty{14.0}{Pa}$, contra $\qty{19.6}{Pa}$ para um canal muito largo de \qty{2}{m} de profundidade: a razão é $R_h/h = 0.714$, as margens tomando o resto do atrito. (c)~Para um canal retangular de largura $b$, $R_h = bh/(b + 2h) = h/(1 + 2h/b)$, que tende a $h$ quando $b$ cresce.
\end{solution}

\begin{solution}{exo:b1:erosion-transport:6}
(a)~$d_1 = \left(18\nu^2/(Rg)\right)^{1/3} = (18/(Rg))^{1/3}\,\nu^{2/3}$, e $R$ e $g$ não dependem da viscosidade. (b)~$\nu = \num{1.3059e-3}/999.7 =
\qty{1.306e-6}{m^2.s^{-1}}$ a \qty{10}{\celsius}, contra
\qty{1.0034e-6}{m^2.s^{-1}} a \qty{20}{\celsius}: a razão é \num{1.302},
e $1.302^{2/3} = 1.192$, logo $d_1 = 0.104 \times 1.192 = \qty{0.124}{mm}$. (c)~No intervalo de Stokes $w_s$ é inversamente proporcional a $\nu$: o grão afunda $1.302$ vezes mais devagar, cerca de \qty{30}{\percent}. Água fria mantém o sedimento fino em suspensão por mais tempo: um rio no inverno carrega seu silte um pouco mais longe, para a mesma vazão.
\end{solution}

\begin{solution}{exo:b1:erosion-transport:7}
Com $R g/(18\nu) = \qty{8.99e5}{m^{-1}.s^{-1}}$ e a lei combinada:
\begin{center}
\begin{tabular}{@{}lccc@{}}
\hline
$d$ & Stokes (\si{m.s^{-1}}) & combinada (\si{m.s^{-1}}) & razão\\
\hline
\qty{0.1}{mm} & \num{8.99e-3} & \num{7.53e-3} & \num{1.19}\\
\qty{0.5}{mm} & \num{0.225} & \num{0.0711} & \num{3.16}\\
\qty{2}{mm} & \num{3.60} & \num{0.197} & \num{18.3}\\
\hline
\end{tabular}
\end{center}
A razão cresce com o tamanho, porque o número de Reynolds do grão cresce: a lei de Stokes, válida apenas para números de Reynolds pequenos, cresce como $d^2$, enquanto o arrasto da esteira freia os grãos grossos até uma velocidade que cresce apenas como $\sqrt d$.
\end{solution}

\begin{solution}{exo:b1:erosion-transport:8}
(a)~$(Rg/\nu^2)^{1/3} = \qty{25263}{m^{-1}}$. Em \qty{0.2}{mm}, $D_* = 5.05$
e $\theta_c = 0.3/7.06 + 0.055(1 - e^{-0.101}) = 0.0425 + 0.0053 =
0.0478$; $\tau_c = 0.0478 \times 16204 \times \num{2e-4} = \qty{0.155}{Pa}$.
Em \qty{20}{mm}, $D_* = 505$ e $\theta_c = 0.3/607 + 0.055(1 - e^{-10.1}) =
0.0005 + 0.0550 = 0.0555$; $\tau_c = 0.0555 \times 16204 \times 0.02 =
\qty{18.0}{Pa}$. (b)~$h = \tau_c/(\rho g S)$ com $\rho g S =
\qty{19.58}{Pa.m^{-1}}$: \qty{7.9}{mm} de água movem a areia fina, \qty{0.92}{m}
o cascalho.
\end{solution}

\begin{solution}{exo:b1:erosion-transport:9}
A bacia perde \qty{0.8}{Mt} = \qty{8e5}{t} por ano, ou $\num{8e5}/2000 =
\qty{400}{t.km^{-2}.yr^{-1}}$. Pela \cref{prop:b1:erosion-transport:denudation},
$D = \qty{8e8}{kg}/(2650 \times \qty{2e9}{m^2}) = \qty{1.51e-4}{m}$ por ano, ou
\qty{151}{m/Myr}. Isso é cerca de três vezes a mediana das bacias medidas com berílio-10, \qty{54}{m/Myr}, embora abaixo de sua média de
\qty{218}{m/Myr}: uma bacia razoavelmente rápida, como as de países íngremes ou tectonicamente ativos.
\end{solution}

\begin{solution}{exo:b1:erosion-transport:10}
(a)~$u_* = \sqrt{g h S} = \sqrt{9.81 \times 3 \times \num{5e-4}} =
\qty{0.1213}{m.s^{-1}}$ e $\tau_b = \rho u_*^2 = \qty{14.7}{Pa}$.
(b)~A suspensão exige $w_s < u_* = \qty{0.1213}{m.s^{-1}}$. A lei combinada dá \qty{0.126}{m.s^{-1}} a \qty{1}{mm}; resolvendo
$w_s(d) = 0.1213$ por tentativas dá $d = \qty{0.94}{mm}$. (c)~Para o cascalho
$\theta_c \approx 0.0555$, logo $\tau_c = \qty{14.7}{Pa}$ para $d \approx
14.7/(0.0555 \times 16204) = \qty{16.3}{mm}$; nesse tamanho a curva de limiar dá
$\theta_c = 0.0555$, de fato. (d)~Grãos mais grossos que cerca de \qty{16}{mm} permanecem no leito e formam um piso de cascalho estável; o cascalho mais fino e a areia mais grossa que cerca de um milímetro viajam como \omterm{def:b1:erosion-transport:load}{carga de fundo}; areia mais fina, silte e argila viajam em suspensão.
\end{solution}

\begin{solution}{exo:b1:erosion-transport:11}
(a)~$0.30 \times \qty{1000}{m} = \qty{300}{m}$ de rocha, média sobre a
região. (b)~A recuperação é $300 \times 2700/3300 = \qty{245}{m}$. (c)~A superfície média é rebaixada apenas em
$300 - 245 = \qty{55}{m}$. (d)~Os cumes, que não foram erodidos, são levados para cima com toda a região: agora estão
\qty{245}{m} mais altos que antes. Cortar vales pode elevar os picos de uma cordilheira.
\end{solution}

\begin{solution}{exo:b1:erosion-transport:12}
(a)~$0.9 \times \qty{5e8}{kg}/\qty{1300}{kg.m^{-3}} = \qty{346000}{m^3}$ por ano,
ou \qty{0.69}{\percent} do reservatório. (b)~Metade do reservatório,
\qty{25e6}{m^3}, é preenchida em $\num{25e6}/\num{346000} = \num{72}\text{ anos}$ nesse ritmo. (c)~À medida que o reservatório se enche, seu volume, e portanto o tempo de residência da
água, diminui, e grãos mais finos passam por ele: sua eficiência de retenção cai.
Ao mesmo tempo, os depósitos antigos são soterrados e compactados, expulsando água, de modo que
sua densidade aparente aumenta e cada tonelada ocupa menos espaço. Os dois efeitos desaceleram
a perda de volume à medida que o reservatório envelhece.
\end{solution}

\begin{solution}{pb:b1:erosion-transport:1}
\textbf{1.} Um ano de água é $1.2 \times \num{3.156e7} = \qty{3.79e7}{m^3}$;
espalhado sobre \qty{8e7}{m^2}, uma camada de \qty{0.473}{m}: uma vazão de
\qty{473}{mm} por ano.

\textbf{2.} $1.2 \times \qty{0.150}{kg.m^{-3}} = \qty{0.18}{kg.s^{-1}}$, ou
\qty{5680}{t} por ano.

\textbf{3.} $1.2 \times 0.080 = \qty{0.096}{kg.s^{-1}}$, ou \qty{3030}{t} por
ano.

\textbf{4.} A \omterm{def:b1:erosion-transport:load}{carga de fundo} é \qty{568}{t} por ano, e o total $5680 + 3030 +
568 = \qty{9278}{t}$ por ano.

\textbf{5.} A concentração de sedimento em suspensão é alta durante os poucos dias de enchente, quando a vazão também é alta, e baixa no resto do tempo. Uma média no tempo dá a cada dia o mesmo peso, de modo que ela é dominada pelos muitos dias de vazão baixa e concentração baixa. Pela
\cref{prop:b1:erosion-transport:flux}, a massa anual é $\int Q\,C\,\mathrm dt$; a média ponderada pela vazão é essa integral dividida por
$\int Q\,\mathrm dt$, de modo que é a média ponderada pela vazão que deve ser multiplicada pela vazão média.

\textbf{6.} $3030/9278 = \num{0.33}$, um terço. Para os rios do mundo, $3.9/(12.6 + 3.9) = \num{0.24}$, e menos se a \omterm{def:b1:erosion-transport:load}{carga de fundo} fosse contada. A parcela da bacia é maior: colinas baixas, onde a erosão mecânica lenta deixa à rocha tempo para se dissolver, e calcário e giz, que se dissolvem por completo.

\textbf{7.} $9278/80 = \qty{116}{t.km^{-2}.yr^{-1}}$.

\textbf{8.} $D = \qty{9.278e6}{kg}/(2700 \times \qty{8e7}{m^2}) =
\qty{4.30e-5}{m}$ por ano, \qty{43.0}{m/Myr}.

\textbf{9.} $43.0/54 = 0.80$: um pouco abaixo da mediana das bacias medidas, e cinco vezes abaixo de sua média. Uma bacia moderadamente lenta, como se espera para colinas baixas em uma região calma.

\textbf{10.} A medição de vazão faz a média sobre os poucos anos de medição, e depende das enchentes que ocorreram neles. O berílio-10 faz a média sobre o tempo necessário para erodir os poucos metros em que ele é produzido: a \qty{43}{m/Myr}, cada metro leva \num{23000} anos, de modo que a média abrange dezenas de milhares de anos. Elas podem diferir por causa de grandes enchentes raras, ausentes em um registro curto, porque a agricultura e o desmatamento aumentaram a erosão nos últimos séculos, e porque o sedimento é armazenado na planície de inundação, e dela retomado.

\textbf{11.} Parte da \omterm{def:b1:erosion-transport:load}{carga dissolvida} não vem da rocha: boa parte do bicarbonato carrega carbono tirado do ar ou do solo, e a chuva traz sal marinho. Contá-los creditaria à bacia uma rocha que ela não perdeu.

\textbf{12.} $\num{1.5e6}/1.2 = \qty{1.25e6}{s}$, \num{14.5} dias.

\textbf{13.} A lei combinada dá \qty{7.53e-3}{m.s^{-1}}: o grão afunda \qty{8}{m} em \qty{1062}{s}, cerca de \qty{18}{min}. Ele é retido.

\textbf{14.} $w_s = \num{8.99e5} \times (\num{2e-5})^2 =
\qty{3.60e-4}{m.s^{-1}}$: \qty{8}{m} em \qty{22200}{s}, cerca de \num{6.2}~horas. Retido.

\textbf{15.} $w_s = \num{8.99e5} \times (\num{1e-6})^2 =
\qty{8.99e-7}{m.s^{-1}}$: \qty{8}{m} leva \qty{8.9e6}{s}, \num{103}~dias, sete vezes o tempo de residência. A maior parte dela passa pela barragem com a água.

\textbf{16.} O grão deve afundar \qty{8}{m} em \qty{1.25e6}{s}, a
$\qty{6.4e-6}{m.s^{-1}}$ no mínimo. Pela lei de Stokes $d = \sqrt{18\nu
w_s/(Rg)} = \sqrt{\num{6.4e-6}/\num{8.99e5}} = \qty{2.7}{\micro m}$: uma argila fina. Grãos mais grossos ficam retidos onde entram.

\textbf{17.} $568 + 0.7 \times 5680 = \qty{4544}{t}$ por ano, o que preenche
$\num{4.544e6}/1300 = \qty{3496}{m^3}$ por ano, \qty{0.233}{\percent} do
reservatório. Metade dele, \qty{7.5e5}{m^3}, é preenchida em cerca de \num{215} anos.

\textbf{18.} $100/43.0 = \qty{2.33}{Myr}$.

\textbf{19.} $1 - 2700/3300 = 0.18$: a superfície desce apenas
\qty{18}{\percent} da espessura removida.

\textbf{20.} $100/0.1818 = \qty{550}{m}$ de rocha.

\textbf{21.} A região inteira se eleva por $550 \times 2700/3300 = \qty{450}{m}$
de recuperação. Cumes pouco erodidos são levados para cima com ela, em até
\qty{450}{m}, enquanto os vales se aprofundam: a topografia cresce antes de declinar.

\textbf{22.} A \omterm{def:b1:erosion-transport:denudation}{taxa de denudação} depende sobretudo das encostas. À medida que as colinas são rebaixadas e suas encostas se suavizam, a taxa diminuirá, a menos que o soerguimento ou uma mudança de clima as torne íngremes outra vez: o tempo calculado à taxa de hoje é um limite inferior.

\textbf{23.} $550/42.95 = 12.80$. À taxa de hoje e com a recuperação isostática, rebaixar as Colinas de Varden em \qty{100}{m} leva \textbf{12.8}~milhões de anos.
\end{solution}
